\subsubsection{Weight Decay Results}
\label{sec:wd}

% Owner: Member 2. Target 1.5--2 pages.

\paragraph{Experiment focus.}
This subsection reports the results of the weight-decay sweep defined in
\cref{sec:method:wd}, evaluated under the shared protocol of
\cref{sec:exp:training,sec:exp:eval}. The only variable that changes is
the AdamW decoupled weight-decay coefficient $\lambda$.

\paragraph{Aggregate numbers.}
\Cref{tab:wd} reports best validation accuracy, final test accuracy,
the train--val gap at the best-val epoch, and wall-clock training
time for each weight-decay setting.

% Weight-decay sweep results. Owner: Member 2.
% Replace -- placeholders with real numbers from the runs.
% Convention: best value per column in \textbf{...}; consistent decimals.

\begin{table}[t]
  \centering
  \small
  \caption{Effect of AdamW weight-decay strength $\lambda$ on
  ViT-Tiny/16 fine-tuned on CIFAR-100. All other hyperparameters fixed
  per \cref{sec:exp:training}. Best per column in bold.}
  \label{tab:wd}
  \begin{tabular}{lcccc}
    \toprule
    Weight decay $\lambda$ & Best val acc (\%) $\uparrow$
                          & Test acc (\%) $\uparrow$
                          & Train--val gap (pp) $\downarrow$
                          & Train time (min) \\
    \midrule
    $0$                   & 84.50 & \textbf{84.75} & 5.38 & 31.5 \\
    $10^{-4}$             & 84.38 & 84.31 & 5.28 & 31.5 \\
    $10^{-2}$             & 84.28 & 84.74 & 5.54 & \textbf{31.3} \\
    $5 \times 10^{-2}$    & 83.88 & 84.05 & 5.54 & 31.5 \\
    $10^{-1}$             & \textbf{84.56} & 84.33 & \textbf{4.97} & 31.7 \\
    \bottomrule
  \end{tabular}
\end{table}


\paragraph{Learning curves.}
\Cref{fig:wd-curves} overlays training and validation loss and
accuracy across all five settings, sharing axes with
\cref{fig:baseline-curves} for direct comparison.

\begin{figure}[t]
  \centering
  % \includegraphics[width=0.9\linewidth]{wd_curves.png}  % uncomment when figure is ready
  \fbox{\rule{0pt}{4cm}\rule{0.85\linewidth}{0pt}}
  \caption{Training and validation curves across weight-decay
  strengths $\lambda \in \{0, 10^{-4}, 10^{-2}, 5\times 10^{-2},
  10^{-1}\}$. \draftnote{Drop \texttt{figures/wd\_curves.png} and
  replace the \texttt{\textbackslash fbox} placeholder with the
  commented \texttt{\textbackslash includegraphics} above. One
  figure, four subplots (train loss, val loss, train acc, val acc),
  one line per $\lambda$, color-coded.}}
  \label{fig:wd-curves}
\end{figure}

\paragraph{Train--val gap.}
\Cref{fig:wd-gap} plots the train--val accuracy gap as a function of
training epoch for each $\lambda$. A flat or shrinking gap at large
$\lambda$ is direct evidence that weight decay is constraining
memorization rather than just trading capacity for regularization.

\begin{figure}[t]
  \centering
  % \includegraphics[width=0.7\linewidth]{wd_gap.png}  % uncomment when figure is ready
  \fbox{\rule{0pt}{4cm}\rule{0.65\linewidth}{0pt}}
  \caption{Train--val accuracy gap over training for each weight-decay
  strength. \draftnote{Drop \texttt{figures/wd\_gap.png} and swap
  in the commented \texttt{\textbackslash includegraphics} above. The
  caption should state the punch line: e.g. ``$\lambda=10^{-2}$
  reduces the final-epoch gap by X percentage points relative to
  $\lambda=0$.''}}
  \label{fig:wd-gap}
\end{figure}

\paragraph{Best validation accuracy as a function of $\lambda$.}
\Cref{fig:wd-best-val} summarizes the sweep by plotting best
validation accuracy against $\lambda$ on a log-x axis. This is the
single plot that answers ``what coefficient should I pick?''

\begin{figure}[t]
  \centering
  % \includegraphics[width=0.6\linewidth]{wd_best_val_vs_wd.png}  % uncomment when ready
  \fbox{\rule{0pt}{4cm}\rule{0.55\linewidth}{0pt}}
  \caption{Best validation accuracy versus weight-decay strength
  $\lambda$ (log scale on the x-axis; $\lambda=0$ is shown as a
  separate marker on the left). \draftnote{Drop
  \texttt{figures/wd\_best\_val\_vs\_wd.png} and swap in the commented
  \texttt{\textbackslash includegraphics} above. Mark the optimum.}}
  \label{fig:wd-best-val}
\end{figure}

\paragraph{Is weight decay effective for fine-tuning ViT-Tiny on
CIFAR-100?}
\label{sec:wd:discussion}
\draftnote{Topic sentence: state whether the best $\lambda > 0$
beats $\lambda = 0$ on test accuracy, and by how many points.
Then 2--3 sentences of evidence: cite the relevant rows of
\cref{tab:wd} and the corresponding curve in \cref{fig:wd-curves}.}

\paragraph{What is roughly the optimal coefficient?}
\draftnote{Topic sentence: name the optimum (likely $10^{-2}$ or
$5\times 10^{-2}$ for AdamW on a small ViT). Then describe the shape
of the curve in \cref{fig:wd-best-val}: is it sharply peaked or a
broad plateau? A broad plateau is the more useful finding because it
implies the practitioner does not need fine tuning.}

\paragraph{Does it act as overfitting prevention or capacity
suppression?}
\draftnote{Use the train--val gap evidence in \cref{fig:wd-gap}.
The decision rule: if increasing $\lambda$ \emph{shrinks} the gap
\emph{while} val accuracy improves or holds, weight decay is
preventing overfitting. If the gap shrinks because train accuracy
collapses faster than val, then weight decay is suppressing capacity
(underfitting). For our largest $\lambda = 10^{-1}$, characterize
which of these two regimes we are in; this is the key qualitative
finding of the section.}

\paragraph{Limitations.}
\label{sec:wd:limitations}
\draftnote{1 short paragraph. Cover: (i) single seed per
configuration, so reported differences may sit within run-to-run
noise---report the magnitude if multi-seed runs are available;
(ii) we study fine-tuning from an ImageNet-pretrained checkpoint,
not training from scratch, so conclusions may not transfer to the
from-scratch regime where weight decay matters more
\citep{touvron2021deit}; (iii) absolute differences between
neighboring $\lambda$ values may be small (sub-1\%), so we are
careful not to over-interpret rank orderings.}
